\subsubsection{Controlled-Controlled-NOT (CCNOT) Gate (Toffoli)}

The \definition{Controlled-Controlled-NOT (CCNOT) Gate}, also called the \definition{Toffoli Gate}, is a \textbf{three-qubit gate} composed of:
\begin{itemize}
    \item Two \textbf{control qubits} that determine whether the gate will be applied.
    \item One \textbf{target qubit} that is flipped (NOT operation) if both control qubits are in the state $\ket{1}$.
\end{itemize}
It \hl{generalizes the CNOT gate by adding an additional control qubit}, making it a crucial component in quantum computing for implementing more complex operations and algorithms.

\highspace
\begin{flushleft}
    \textcolor{Green3}{\faIcon{book} \textbf{Action on Basis States}}
\end{flushleft}
The CCNOT gate acts on the basis states as follows:

\begin{table}[!htp]
    \centering
    \begin{tabular}{@{} c | c @{}}
        \toprule
        \textbf{Input} & \textbf{Output} \\
        \midrule
        $\ket{000}$ & $\ket{000}$ \\
        $\ket{001}$ & $\ket{001}$ \\
        $\ket{010}$ & $\ket{010}$ \\
        $\ket{011}$ & $\ket{011}$ \\
        $\ket{100}$ & $\ket{100}$ \\
        $\ket{101}$ & $\ket{101}$ \\
        $\ket{110}$ & $\ket{111}$ \\
        $\ket{111}$ & $\ket{110}$ \\
        \bottomrule
    \end{tabular}
\end{table}

\noindent
The target flips \textbf{only if both control qubits are $\ket{1}$}, demonstrating the conditional nature of the gate (last two rows of the table). For all other input combinations, the output remains unchanged.

\highspace
\begin{flushleft}
    \textcolor{Green3}{\faIcon{square-root-alt} \textbf{Matrix Representation}}
\end{flushleft}
The matrix representation of the CCNOT gate in the computational basis is an $8 \times 8$ (since it operates on three qubits) unitary matrix given by:
\begin{equation}
    \text{CCNOT} = \begin{bmatrix}
        1 & 0 & 0 & 0 & 0 & 0 & 0 & 0 \\
        0 & 1 & 0 & 0 & 0 & 0 & 0 & 0 \\
        0 & 0 & 1 & 0 & 0 & 0 & 0 & 0 \\
        0 & 0 & 0 & 1 & 0 & 0 & 0 & 0 \\
        0 & 0 & 0 & 0 & 1 & 0 & 0 & 0 \\
        0 & 0 & 0 & 0 & 0 & 1 & 0 & 0 \\
        0 & 0 & 0 & 0 & 0 & 0 & 0 & 1 \\
        0 & 0 & 0 & 0 & 0 & 0 & 1 & 0 \\
    \end{bmatrix}
\end{equation}
This matrix shows that the first six basis states are unchanged (diagonal entries of 1), while the \textbf{last two basis states are swapped} (off-diagonal entries of 1 in the last two rows and columns).

\highspace
\begin{flushleft}
    \textcolor{Green3}{\faIcon{book} \textbf{Action on General States}}
\end{flushleft}
For a general three-qubit state:
\begin{equation*}
    \ket{\psi} = \sum_{i,j,k \in \{0,1\}} \alpha_{ijk} \ket{ijk}
\end{equation*}
The CCNOT gate acts linearly on each component of the superposition and flips the target qubit (the third qubit) if and only if both control qubits are equal to $\ket{1}$. This can be expressed as:
\begin{equation}\label{eq:ccnot-action-general}
    \text{CCNOT} \ket{\psi}
    =
    \sum_{i,j,k \in \{0,1\}} \alpha_{ijk} \ket{ij\left(k \oplus (i \cdot j)\right)}
\end{equation}
Where $\oplus$ denotes the XOR operation, and $i \cdot j = 1$ if and only if $i = j = 1$, and $0$ otherwise. 

\highspace
In simple terms, applying \hl{CCNOT affects only the $c_6$ and $c_7$ coefficients} (corresponding to $\ket{110}$ and $\ket{111}$) by swapping them, while all \hl{other coefficients remain unchanged.}

\highspace
\begin{flushleft}
    \textcolor{Green3}{\faIcon{tools} \textbf{Circuit Representation}}
\end{flushleft}
In quantum circuit diagrams, the CCNOT gate is typically represented by a symbol with two control dots connected to a NOT gate (X gate) on the target qubit. The control dots indicate which qubits are the controls, and the NOT gate indicates the operation performed on the target qubit when both controls are active. It is similar to the CNOT gate but with an additional control dot, making it visually distinct and easy to identify in quantum circuits.

\begin{center}
    \begin{quantikz}
        \lstick{$\ket{v_0}$} & \ctrl{1} & \rstick{$\ket{v_0}$} \\
        \lstick{$\ket{v_1}$} & \ctrl{1} & \rstick{$\ket{v_1}$} \\
        \lstick{$\ket{v_2}$} & \targ{}  & \rstick{$\ket{v_2} \oplus \left(\ket{v_0} \land \ket{v_1}\right)$} \\
    \end{quantikz}
\end{center}

\noindent
In summary, \hl{if both control qubits are $\ket{1}$, apply NOT to the target; otherwise do nothing.}

\highspace
\begin{flushleft}
    \textcolor{Green3}{\faIcon{cogs} \textbf{Properties}}
\end{flushleft}
\begin{itemize}
    \item \important{Unitary}
    \begin{equation*}
        \text{CCNOT}^\dagger \text{CCNOT} = I
    \end{equation*}
    
    
    \item \important{Reversible}
    \begin{equation*}
        \text{CCNOT}^{-1} = \text{CCNOT}^\dagger
    \end{equation*}
    

    \item \important{Classical universal gate (can simulate AND)}. CCNOT gate is a \textbf{universal gate} for classical computation, meaning that \textbf{any classical logic gate can be constructed using CCNOT gates and single-qubit NOT gates}. In particular, it can compute the AND operation in a reversible way:
    \begin{equation*}
        \ket{a,\, b,\, 0} \xrightarrow{\text{CCNOT}} \ket{a,\, b,\, 0 \oplus \left(a \cdot b\right)} = \ket{a,\, b,\, a \cdot b}
    \end{equation*}
    Where $a$ and $b$ are the control qubits, and the target qubit is initialized to $\ket{0}$. After applying the CCNOT gate, the target qubit will be in the state $\ket{1}$ if and only if both $a$ and $b$ are $\ket{1}$, effectively computing the AND of the two control qubits.

    Since AND and NOT are sufficient to express any Boolean function, the Toffoli gate can reproduce any classical computation.
    

    \item \important{Introduce multi-qubit interaction}. Multi-qubit interaction refers to operations in which the \hl{evolution of one qubit depends on the state of one or more other qubits.} Unlike parallel gates, which act independently on each qubit, \textbf{interacting gates introduce conditional behavior and correlations between qubits}.

    Typical examples include controlled gates such as CNOT and CCNOT, where the transformation applied to a target qubit depends on the value of one or more control qubits. These interactions are essential for creating entanglement and for implementing complex quantum algorithms.
    

    \item \important{Can create entanglement}. Similar to the CNOT gate, also the CCNOT gate can create entanglement between the control and target qubits. For example, if the control qubits are in a superposition state, applying the CCNOT gate can result in an entangled state between the control and target qubits.
    

    \item \important{Not decomposable into a single tensor product}. An interacting multi-qubit gate cannot be expressed as a tensor product of single-qubit operators. That is, in general:
    \begin{equation*}
        U \neq U_A \otimes U_B
    \end{equation*}
    This means that the gate introduces correlations between qubits and cannot be reduced to independent operations.
\end{itemize}
In summary, the Toffoli gate performs a conditional operation based on two qubits, enabling multi-qubit logic and more complex correlations.